\begin{proof}

\begin{enumerate}[label={$\langle1\rangle$\arabic*.},leftmargin=*]
\item Use $\rho:T\to N\to S$ and the effective
integral divisor $D$ of P4R.6. Put
$A=\operatorname{Supp}D$ and
$C=\operatorname{Sing}(S)_{\rm red}$.
The Hopf-curve and blowup calculations of
P4R.5 use only integrality, effectiveness,
and $K_T+D\equiv0$. They therefore apply
to this divisor: $A$ has $r\ge1$ connected
components, each a smooth elliptic root
with rational trees attached. Each of its
$v$ rational components came from one of
the $k$ blowups; $e(A)=v$, $e(T)=k$, and
$\rho$ is an isomorphism off $A$. This
proves the displayed formula for the
regular locus.

\item Every one-dimensional component of $C$
receives a nonconstant proper map from a
component of $A$, so its normalization
has genus zero or one. There are no
isolated points of $C$: an isolated
singular point of the hypersurface is
normal by the hypersurface $S_2$ property
and regularity in codimension one, and
P4R.6's vanishing-sphere argument excludes
it. Let $s,c,g,b$ respectively count the
connected components of $C$, its
irreducible components, its genus-one
normalizations, and the cycle ranks of
its point-identification graphs.
Then $e(C)=c+s-2g-b$, $r\ge g$, and
$c\ge g$.

\item The proper-map pushout $T\cup_A C=S$
and its cofibration justification are
exactly those of P4R.7. A root which
maps to a point gives the zero
homomorphism on its fundamental group;
a noncontracted root maps through the
normalization of one component of $C$.
Every rational tree is simply connected.
Consequently killing $\pi_1(T)$ and
the normalized component groups of $C$
in the disconnected van Kampen
presentation leaves the free quotient
$F_{b+r-s}$, as in P4R.5. Virtual
abelianness gives $b+r-s\le1$, hence
\[
 e(C)\ge c+r-2g-1\ge-1.
\]

\item A negative value forces $c=r=g$.
Every root is then noncontracted, maps
onto an elliptic normalization, and
supplies one of the $c$ image components.
Every rational component of $A$ is
contracted: a nonconstant map from its
normalization $\mathbb P^1$ to an
elliptic curve contradicts
Riemann--Hurwitz. Thus every exceptional
block is a rational subtree of one
rooted component of $A$. It meets that
root in exactly one point. At least one
intersection is forced by the final
assertion of P4R.6, since the roots are
the only strict conductor components;
two intersections would make a cycle
in the rooted tree. The exceptional
fibre is connected. One can see this
directly by Stein factorization of
$T\to N$: its finite birational factor
is an isomorphism because $N$ is normal.

\item The conductor on $N$ consequently has
one local analytic branch at each of
its points. At a singular point of $N$,
each local branch has a strict transform
meeting the exceptional fibre, whereas
there is exactly one such intersection
with the smooth root. Away from those
fibres the assertion follows from the
smooth disjoint roots themselves.

\item The local hypersurface argument of
P4R.5 now applies with $N$ in place of
its smooth resolution. Two different
curve germs of $C$ would require two
different normalization points, because
each point has just one conductor germ.
Their corresponding hypersurface
branches intersect in a curve germ.
That germ belongs to both conductor
images and forces the two chosen germs
to coincide. Thus $C$ is locally
unibranch with disjoint irreducible
components. This gives $b=0$, $s=c=g$,
and $e(C)=0$, a contradiction.
Therefore $e(C)\ge0$, proving the
assertion.

\end{enumerate}
\end{proof}
